%!TEX root = TheoretischePhysikC.tex

\section{Vorlesung: Das Coulombpotential (Wasserstoffatom)}

\begin{center}
	\begin{tikzpicture}
		\begin{axis}[width=\columnwidth, height=8cm,hide axis, tick style={color=black}, domain=0:12, xtick=\empty, ytick=\empty, samples =200, xmin = -0.5, xmax = 8, ymin = -3, ymax = 2]
		\addplot[cgblue, dashed] {1/x^2-2/x};
		\addplot[cgblue] {-2/x};
		\draw[->] (axis cs: -0.5, 0) -- (axis cs: 8, 0);
		\draw (axis cs: 7.8, -0.2) node[below] {$r$};
		\draw (axis cs: 6, 1.5) node[below] {$U_\ell(r) = \frac{\hbar^2 \ell(\ell+1)}{2m_{\text{e}}r^2}+V(r)$};
		\draw[->] (axis cs: 0, -3) -- (axis cs: 0, 2);
    	\draw (axis cs: 0.5,1.5) node[right] {$\ell = 1$};
   		\draw (axis cs:3.5,-1.25) node[above] {$\ell=0$};
		\end{axis}
	\end{tikzpicture}
\end{center} 

Das Wasserstoffatom besteht aus einem Proton und einem Elektron. Ein erstes grobes Modell für das Atom erfolgt indem man das Proton als eine unendlichen massiven Punktmasse annähert. Nach dieser Annäherung kann nun das Wasserstoffatom als ein Teilchen (das Elektron), das sich in einem Zentralpotential 
\begin{equation}
	V(r) = -\frac{Ze^2}{4\pi \epsilon_0 r}
\end{equation}
bewegen, modelliert werden.

Das Elektron ist dabei kein Punkt auf einer Bahn, sondern eine Wolke; das Proton
ebenso, nur lässt der Massenunterschied es als Punktmasse durchgehen. Die Reduktion
auf ein Teilchen ist keine Näherung: Relativkoordinaten und die reduzierte Masse
$\mu = m_{\text{e}}m_{\text{p}}/(m_{\text{e}}+m_{\text{p}})$ machen sie exakt;
genähert wird erst mit $\mu \approx m_{\text{e}}$. Aus solchen Rechnungen wurden
die Wasserstofflinien bestimmt, mit denen man Entfernungen zu fernen Galaxien
vermaß.

Die Radialgleichung ist ein eindimensionales Problem im effektiven Potential
$U_\ell(r)$ des Bildes oben. Lösen lässt sie sich mit dem Potenzreihenansatz, der
Frobenius-Methode; das ist so aufwendig, dass wir die Lösungen einfach
hinschreiben und durch Einsetzen prüfen.

Die Lösungen der Radialgleichung 
\begin{equation}
	\left(-\frac{\hbar^2}{2m_{\text{e}}}\frac{d^2}{dr^2} + \frac{\hbar^2\ell(\ell+1)}{2m_{\text{e}}r^2} + V(r) - E\right)u(r) = 0.
\end{equation}
sind der Form
\begin{equation}
	R_{n\ell}(r) = \sqrt{\frac{(n-\ell-1)!(2\kappa)^3}{2n(n+\ell)!}}(2\kappa r)^\ell e^{-\kappa r}\text{L}_{n-\ell-1}^{2\ell +1}(2\kappa r),	
\end{equation}
wobei 
\begin{equation}
	\kappa = \frac{\sqrt{2m_{\text{e}}|E|}}{\hbar} = \frac{Z}{n a_0}
\end{equation}
und $L_r^s$ die Laguerre-Polynomen sind. Die dazugehörige Eigenwerte lauten
\begin{equation}
	\eqbox{E_n = -\frac{(Ze)^2}{4\pi\epsilon_0}\frac{1}{a_0}\frac{1}{2n^2},}
\end{equation}
mit \emph{Hauptquantenzahl} $n=1,2,\ldots$ und \emph{Bohrradius} 
\begin{equation}
	a_0 = \frac{4\pi\epsilon_0\hbar^2}{m_{\text{e}}e^2} = 5.29177210903(80) \times 10^{-11} \text{m}.
\end{equation}
Die Zählung beginnt bei $n=1$, ein $n=0$ gibt es nicht; im Termschema häufen sich
die unendlich vielen Niveaus an der Ionisationsschwelle $E=0$.
\begin{merksatz}
Beobachtbar sind nicht die Eigenwerte, sondern ihre Differenzen: beim Übergang
$n \to n'$ mit $n'<n$ emittiert das Elektron ein Photon der Energie $E_n-E_{n'}$,
und das sind die Spektrallinien.
\end{merksatz}

% fig_vl16_termschema.tex — Das Termschema des Wasserstoffs und die
% Spektrallinien, die aus seinen Niveaudifferenzen folgen.
%
% Follows the conventions of fig_vl11_kasten (the reference plate).
%
%   * geometry comes from the closed form derived in the lecture, never
%     from a freehand sketch.  With the Rydberg energy
%         E_R = e^2/(4 pi eps_0 a_0) / 2 = 13.6057 eV,
%     the lecture's eigenvalue (Z = 1)
%         E_n = -E_R / n^2
%     fixes every level height in panel (a):
%         n  = 1      2       3       4       5       6      ...
%         E  = -13.61 -3.401 -1.512  -0.850  -0.544  -0.378  ... -> 0
%     so the 1/n^2 accumulation at the ionisation threshold is forced by
%     the formula and not drawn by eye.  Levels up to n = 20 are drawn;
%     from n ~ 8 they can no longer be resolved on the plate, which is
%     precisely the statement "infinitely many levels pile up at E = 0".
%   * the transition arrows carry the Merksatz of the lecture: observable
%     are not the eigenvalues but their differences.  Hence the levels are
%     neutral grey reference rules and the *arrows* are the coloured data.
%     Arrow length is literally the photon energy E_n - E_n', to scale:
%     the Lyman arrows are ~4x the Balmer arrows, which is the whole
%     reason the one series is deep UV and the other visible light.
%   * colour encodes the final state n' of the series, and nothing else:
%     n' = 1 (Lyman) cgblue, n' = 2 (Balmer) birrenorange — the same two
%     colours in both panels, so a stick in (b) can be traced back to an
%     arrow in (a).
%   * panel (b) plots the wavelengths that follow from the same formula,
%         1/lambda = R_inf (1/n'^2 - 1/n^2),  1/R_inf = 91.127 nm,
%     on a logarithmic axis, because only a log axis holds Lyman
%     (91-122 nm) and Balmer (365-656 nm) in one frame.  Consequences
%     that are therefore forced: the two series never overlap, the band
%     122-365 nm carries no hydrogen line at all, and the lines of each
%     series crowd at their series limit lambda = n'^2 / R_inf exactly as
%     the levels crowd at E = 0.  The visible window 380-750 nm is drawn
%     as a plain grey bracket, not a coloured band: it is reference
%     furniture and must survive greyscale printing.
%
% Panel (a) the term scheme with the Lyman and Balmer transitions.
% Panel (b) the same transitions as spectral lines on a wavelength axis —
%           the order of magnitude prose cannot deliver: Balmer lands in
%           the visible, Lyman does not.

\begin{center}
\begin{tikzpicture}

\tikzset{
  trans/.style={-{Stealth[length=5pt,width=3.5pt]}, line width=0.9pt},
}

% ─────────────────────────────────────────────────────────────
%  (a)  Term scheme E_n = -13.6 eV / n^2
% ─────────────────────────────────────────────────────────────
\begin{axis}[
    name=terme,
    tpcplot,
    width=6.2cm, height=6.9cm,
    axis y line=left, axis x line=none,
    grid=none, clip=false,
    xmin=0, xmax=1.20,
    ymin=-16.4, ymax=1.6,
    ytick={0,-3.4014,-13.6057},
    yticklabels={$0$,$-3.4$,$-13.6$},
    ylabel={$E$ [eV]},
    every axis y label/.style={at={(axis description cs:-0.20,0.5)},
                               rotate=90, anchor=south, font=\small},
    title={(a) \tpcifenglish{Term scheme}{Termschema}},
]

  % Bound levels E_n = -13.6057/n^2.  Written out rather than looped:
  % \foreach defers expansion inside a pgfplots axis.
  \draw[levellinegrey] (axis cs:0.02,-13.6057) -- (axis cs:0.88,-13.6057);
  \draw[levellinegrey] (axis cs:0.02, -3.4014) -- (axis cs:0.88, -3.4014);
  \draw[levellinegrey] (axis cs:0.02, -1.5117) -- (axis cs:0.88, -1.5117);
  \draw[levellinegrey] (axis cs:0.02, -0.8504) -- (axis cs:0.88, -0.8504);
  \draw[levellinegrey] (axis cs:0.02, -0.5442) -- (axis cs:0.88, -0.5442);
  \draw[levellinegrey] (axis cs:0.02, -0.3779) -- (axis cs:0.88, -0.3779);
  \draw[levellinegrey] (axis cs:0.02, -0.2777) -- (axis cs:0.88, -0.2777);
  \draw[levellinegrey] (axis cs:0.02, -0.2126) -- (axis cs:0.88, -0.2126);
  \draw[levellinegrey] (axis cs:0.02, -0.1680) -- (axis cs:0.88, -0.1680);
  \draw[levellinegrey] (axis cs:0.02, -0.1361) -- (axis cs:0.88, -0.1361);
  \draw[levellinegrey] (axis cs:0.02, -0.0945) -- (axis cs:0.88, -0.0945);

  % The ionisation threshold: a real boundary, so heavy ink.  It needs no
  % "E = 0" tag at its right end — the y tick already reads 0 — and that
  % end is where the accumulating levels have to be labelled.
  \draw[potential] (axis cs:0.02,0) -- (axis cs:0.88,0);
  \node[lbl, text=birrengray, font=\scriptsize, anchor=south]
    at (axis cs:0.58,0.10)
    {\tpcifenglish{ionisation threshold}{Ionisationsschwelle}};

  % Lyman series, n -> 1.  Arrow length = E_n - E_1, to scale.
  \draw[trans, cgblue] (axis cs:0.16, -3.4014) -- (axis cs:0.16,-13.6057);
  \draw[trans, cgblue] (axis cs:0.23, -1.5117) -- (axis cs:0.23,-13.6057);
  \draw[trans, cgblue] (axis cs:0.30, -0.8504) -- (axis cs:0.30,-13.6057);

  % Balmer series, n -> 2
  \draw[trans, birrenorange] (axis cs:0.60,-1.5117) -- (axis cs:0.60,-3.4014);
  \draw[trans, birrenorange] (axis cs:0.67,-0.8504) -- (axis cs:0.67,-3.4014);
  \draw[trans, birrenorange] (axis cs:0.74,-0.5442) -- (axis cs:0.74,-3.4014);

  % Series labels, directly under the level each series ends on
  \node[lbl, text=cgblue, anchor=north, align=center]
    at (axis cs:0.30,-13.90) {Lyman ($n\to1$)\\[-1pt] UV};
  \node[lbl, text=birrenorange, anchor=north, align=center]
    at (axis cs:0.70,-3.72)
    {Balmer ($n\to2$)\\[-1pt] \tpcifenglish{visible}{sichtbar}};

  % Quantum numbers, only where the levels are still resolved
  \node[lbl, text=birrengray, font=\scriptsize, anchor=west]
    at (axis cs:0.90,-13.6057) {$n=1$};
  \node[lbl, text=birrengray, font=\scriptsize, anchor=west]
    at (axis cs:0.90, -3.4014) {$n=2$};
  \node[lbl, text=birrengray, font=\scriptsize, anchor=west]
    at (axis cs:0.90, -1.5117) {$n=3$};

  % The accumulation at the threshold: labelled exactly like the resolved
  % levels, in the same column, because that is what it is — the rest of
  % the same sequence, no longer separable on the plate.
  \node[lbl, text=birrengray, font=\scriptsize, anchor=west]
    at (axis cs:0.90,-0.30) {$n\to\infty$};
\end{axis}

% ─────────────────────────────────────────────────────────────
%  (b)  The same transitions as spectral lines
% ─────────────────────────────────────────────────────────────
\begin{axis}[
    tpcplot,
    at={(terme.north east)}, anchor=north west, xshift=13mm,
    width=7.4cm, height=5.7cm,
    axis y line=none, axis x line=bottom,
    grid=none, clip=false,
    xmode=log, log ticks with fixed point,
    xmin=82, xmax=900,
    ymin=0,   ymax=1.42,
    xtick={100,200,400,800},
    minor xtick={90,300,500,600,700},
    xlabel={$\lambda$ [nm]},
    every axis x label/.style={at={(axis description cs:0.5,-0.26)},
                               anchor=north, font=\small},
    title={(b) \tpcifenglish{Spectral lines}{Spektrallinien}},
]

  % Lyman, n = 2..12 -> 1.  lambda = 91.127 nm / (1 - 1/n^2)
  \draw[cgblue, line width=0.7pt] (axis cs:121.50,0) -- (axis cs:121.50,1);
  \draw[cgblue, line width=0.7pt] (axis cs:102.52,0) -- (axis cs:102.52,1);
  \draw[cgblue, line width=0.7pt] (axis cs: 97.20,0) -- (axis cs: 97.20,1);
  \draw[cgblue, line width=0.7pt] (axis cs: 94.92,0) -- (axis cs: 94.92,1);
  \draw[cgblue, line width=0.7pt] (axis cs: 93.73,0) -- (axis cs: 93.73,1);
  \draw[cgblue, line width=0.7pt] (axis cs: 93.03,0) -- (axis cs: 93.03,1);
  \draw[cgblue, line width=0.7pt] (axis cs: 92.57,0) -- (axis cs: 92.57,1);
  \draw[cgblue, line width=0.7pt] (axis cs: 92.27,0) -- (axis cs: 92.27,1);
  \draw[cgblue, line width=0.7pt] (axis cs: 92.05,0) -- (axis cs: 92.05,1);
  \draw[cgblue, line width=0.7pt] (axis cs: 91.76,0) -- (axis cs: 91.76,1);

  % Balmer, n = 3..12 -> 2.  lambda = 91.127 nm / (1/4 - 1/n^2)
  \draw[birrenorange, line width=0.7pt] (axis cs:656.11,0) -- (axis cs:656.11,1);
  \draw[birrenorange, line width=0.7pt] (axis cs:486.01,0) -- (axis cs:486.01,1);
  \draw[birrenorange, line width=0.7pt] (axis cs:433.94,0) -- (axis cs:433.94,1);
  \draw[birrenorange, line width=0.7pt] (axis cs:410.07,0) -- (axis cs:410.07,1);
  \draw[birrenorange, line width=0.7pt] (axis cs:396.91,0) -- (axis cs:396.91,1);
  \draw[birrenorange, line width=0.7pt] (axis cs:388.81,0) -- (axis cs:388.81,1);
  \draw[birrenorange, line width=0.7pt] (axis cs:383.44,0) -- (axis cs:383.44,1);
  \draw[birrenorange, line width=0.7pt] (axis cs:379.69,0) -- (axis cs:379.69,1);
  \draw[birrenorange, line width=0.7pt] (axis cs:374.92,0) -- (axis cs:374.92,1);

  % Series limits lambda = n'^2 / R_inf, where n -> infinity accumulates —
  % the image of the pile-up at E = 0 in panel (a).  Annotated once, from
  % the empty band between the series, so no label sits on a line.
  \draw[levellinegrey] (axis cs: 91.13,0) -- (axis cs: 91.13,1.20);
  \draw[levellinegrey] (axis cs:364.51,0) -- (axis cs:364.51,1.20);
  \node[lbl, text=birrengray, font=\scriptsize, anchor=south west]
    at (axis cs:93.5,1.20) {\tpcifenglish{series limit}{Seriengrenze}};
  \node[lbl, text=birrengray, font=\scriptsize, anchor=south east]
    at (axis cs:358,1.02) {\tpcifenglish{series limit}{Seriengrenze}};

  % The band between the two series is empty, and that emptiness is a
  % result, not a gap in the drawing — hence it is labelled.
  \node[lbl, text=birrengray, font=\scriptsize, align=center]
    at (axis cs:210,0.46)
    {\tpcifenglish{no hydrogen}{keine}\\[-1pt]
     \tpcifenglish{lines here}{Wasserstofflinien}};

  % The one hydrogen line everybody has seen: red H-alpha, n = 3 -> 2.
  \node[lbl, text=birrenorange, font=\scriptsize, anchor=south]
    at (axis cs:656.11,1.00) {$\mathrm{H}\alpha$};

  % Spectral windows, as neutral grey brackets under the axis: reference
  % furniture, so no colour, and greyscale-safe.  The visible window
  % starts just longward of the Balmer limit — which is why that whole
  % series, and only that series, is a laboratory-bench phenomenon.
  \draw[birrengray, line width=0.5pt]
    (axis cs:84,-0.145) -- (axis cs:84,-0.18) -- (axis cs:378,-0.18)
      -- (axis cs:378,-0.145);
  \node[lbl, text=birrengray, font=\scriptsize, anchor=north]
    at (axis cs:178,-0.19) {UV};
  \draw[birrengray, line width=0.5pt]
    (axis cs:384,-0.145) -- (axis cs:384,-0.18) -- (axis cs:750,-0.18)
      -- (axis cs:750,-0.145);
  \node[lbl, text=birrengray, font=\scriptsize, anchor=north]
    at (axis cs:537,-0.19) {\tpcifenglish{visible}{sichtbar}};
\end{axis}

\end{tikzpicture}
\end{center}


Farbig sind hier nicht die Niveaus, sondern die Übergänge: die Pfeillänge ist
maßstäblich die Photonenenergie, die Lyman-Pfeile sind rund viermal so lang wie
die Balmer-Pfeile. Rechts kehrt die Häufung bei $E=0$ als Gedränge an der
Seriengrenze wieder.

Die vollständige Eigenfunktionen lauten
\begin{equation}
	\psi_{n\ell m}(r,\theta,\phi) = R_{n\ell}(r)\text{Y}_{\ell}^m(\theta,\phi).
\end{equation}
Zu jedem $n$ gibt es mögliche Drehimpulsquantenzahlen $\ell = 0,1,2, \ldots, n-1$ ($\equiv \text{s}, \text{p}, \text{d}, \text{f}, \ldots $). Zu jedem $\ell$ gehören $2\ell +1$ mögliche Werte der magnetischen Quantenzahahl $m = -\ell, -\ell+1, \ldots, \ell$. Eine wichtige Bemerkung: die Energie ist von $m$ und $\ell$ unabhängig. 

Die ersten Laguerre-Polynomen lauten
\begin{equation}
	\begin{split}
		L_0^s(x) &= 1 \\
		L_1^s(x) &= -x+(s+1) \\
		L_2^s(x) &= \frac{x^2}{2}-(s+2)x + \frac{1}{2}(s+1)(s+2) \\
		&\vdots
	\end{split}
\end{equation}
Daher erhalten wir die folgende Eigenwellenfunktionen $\psi_{n\ell m}(r,\theta,\phi)$ bis $n=2$:
\begin{align*}
	\psi_{100}(r,\theta,\phi) &= \frac{1}{\sqrt{\pi}}\left(\frac{Z}{a_0}\right)^{\frac{3}{2}}e^{-\frac{Zr}{a_0}} &\quad & 1\text{s} \\
	\psi_{200}(r,\theta,\phi) &= \frac{1}{4\sqrt{2\pi}}\left(\frac{Z}{a_0}\right)^{\frac{3}{2}}\left(2-\frac{Zr}{a_0}\right)e^{-\frac{Zr}{2a_0}} &\quad & 2\text{s} \\
	\psi_{210}(r,\theta,\phi) &= \frac{1}{4\sqrt{2\pi}}\left(\frac{Z}{a_0}\right)^{\frac{3}{2}}\frac{Zr}{a_0}e^{-\frac{Zr}{2a_0}}\cos(\theta) &\quad & 2\text{p} \quad (m=0) \\
	\psi_{21\pm1}(r,\theta,\phi) &= \mp\frac{1}{8\sqrt{\pi}}\left(\frac{Z}{a_0}\right)^{\frac{3}{2}}\frac{Zr}{a_0}e^{-\frac{Zr}{2a_0}}\sin(\theta)e^{\pm i\phi} &\quad & 2\text{p} \quad (m=\pm1). \\
\end{align*}

Die üblichen bunten Orbitalbilder zeigen nur $|\psi_{n\ell m}|^2$, nicht die
Phase.
